% Curso intensivo de los esquemas de demostración de los capítulos 7 a 10 de la
% tesis, en escritura matemática (sin Lean). Fuente: tesis-mate/documento/capitulos
%   1_contextohistorico_matematico_2.tex  (Cap. 7)
%   1_contextohistorico_matematico_3.tex  (Cap. 8)
%   2_linkeo.tex                          (Cap. 9)
%   3_linkeo.tex                          (Cap. 10)
% Compilar dos veces:  lualatex CursoTeoria.tex
\documentclass[11pt,a4paper]{article}
\usepackage[a4paper,margin=2.3cm,headheight=14pt]{geometry}
\usepackage{amsmath,mathtools}
\usepackage{fontspec}
\usepackage{unicode-math}
\usepackage[spanish,es-noshorthands,es-tabla]{babel}
\setmainfont{TeX Gyre Pagella}
\setmathfont{TeX Gyre Pagella Math}
\setsansfont{DejaVu Sans}[Scale=MatchLowercase]
\usepackage{xcolor}
\usepackage[most]{tcolorbox}
\usepackage{booktabs,tabularx,array,longtable}
\usepackage{bussproofs}
\usepackage{flagderiv}
\usepackage{tikz}
\usetikzlibrary{arrows.meta,positioning}
\usepackage{tikz-cd}
\usepackage{enumitem}
\usepackage{fancyhdr}
\usepackage[hidelinks]{hyperref}
\usepackage{microtype}
\AtBeginDocument{\providecommand\square{\mdlgwhtsquare}}

\definecolor{fondoesquema}{HTML}{EEF5EE}
\definecolor{bordeesquema}{HTML}{5E9A62}
\definecolor{fondonota}{HTML}{F1EEF8}
\definecolor{bordenota}{HTML}{8C7DB8}
\definecolor{fondoejemplo}{HTML}{F7F8FA}
\definecolor{bordeejemplo}{HTML}{B8C2CF}
\definecolor{fondotrampa}{HTML}{FFF4EC}
\definecolor{bordetrampa}{HTML}{D98A4E}
\definecolor{okverde}{HTML}{2E7D32}
\definecolor{esbozo}{HTML}{B7791F}
\definecolor{cita}{HTML}{6B6B9A}

\newtcolorbox{esquema}[1]{enhanced, breakable, colback=fondoesquema,
  colframe=bordeesquema, boxrule=0.6pt, arc=3pt,
  colbacktitle=bordeesquema, coltitle=white, fonttitle=\sffamily\bfseries,
  title={#1}, before skip=10pt, after skip=10pt, before upper={\raggedright}}
\newtcolorbox{ejemplotesis}[1]{enhanced, breakable, colback=fondoejemplo,
  colframe=bordeejemplo, boxrule=0.5pt, arc=2pt, colbacktitle=bordeejemplo!40,
  coltitle=black, fonttitle=\sffamily\small\bfseries, title={#1},
  before skip=6pt, after skip=8pt}
\newtcolorbox{nota}[1][Nota]{enhanced, breakable, colback=fondonota,
  colframe=bordenota, boxrule=0.5pt, arc=3pt, colbacktitle=bordenota,
  coltitle=white, fonttitle=\sffamily\bfseries\small, title={#1},
  before skip=8pt, after skip=8pt}
\newtcolorbox{trampa}[1][Lo que hay que vigilar]{enhanced, breakable,
  colback=fondotrampa, colframe=bordetrampa, boxrule=0.5pt, arc=3pt,
  colbacktitle=bordetrampa, coltitle=white, fonttitle=\sffamily\bfseries\small,
  title={#1}, before skip=8pt, after skip=8pt}

\newcommand{\forma}{\par\noindent\textbf{Forma.}\enspace}
\newcommand{\porque}{\par\smallskip\noindent\textbf{Por qué es válido.}\enspace}
\newcommand{\plantilla}{\par\smallskip\noindent\textbf{Cómo se redacta.}\enspace}
\newcommand{\donde}{\par\smallskip\noindent\textbf{Dónde aparece.}\enspace}

\newcommand{\etiqueta}[2]{\tcbox[on line, boxsep=0pt, left=3pt, right=3pt,
  top=1pt, bottom=1pt, colback=#1!12, colframe=#1, boxrule=0.4pt,
  arc=2pt]{\sffamily\scriptsize #2}}
\newcommand{\dem}{\etiqueta{okverde}{demostrado}}
\newcommand{\esb}{\etiqueta{esbozo}{esbozado}}
\newcommand{\cit}{\etiqueta{cita}{citado}}
\newcommand{\enu}{\etiqueta{cita}{enunciado}}

% Notación
\newcommand{\Sub}{\operatorname{Sub}}
\newcommand{\FV}{\operatorname{FV}}
\newcommand{\dom}{\operatorname{dom}}
\newcommand{\bred}{\rightarrow_\beta}
\newcommand{\bmred}{\twoheadrightarrow_\beta}
\newcommand{\beq}{=_\beta}
\newcommand{\aeq}{=_\alpha}
\newcommand{\bconv}{\leftrightarrow_\beta}
\newcommand{\pr}{\Rightarrow}
\newcommand{\sus}[3]{#1[#2:=#3]}
\newcommand{\Elim}{\mathrm{Elim}}
\newcommand{\Ind}{\mathrm{Ind}}
\newcommand{\Constr}{\mathrm{Constr}}

\pagestyle{fancy}
\fancyhf{}
\fancyhead[L]{\small\sffamily Esquemas de demostración de los capítulos 7 a 10}
\fancyhead[R]{\small\sffamily\nouppercase{\leftmark}}
\fancyfoot[C]{\small\thepage}
\renewcommand{\headrulewidth}{0.4pt}

\setlist{itemsep=2pt, topsep=4pt}
\setlength{\parskip}{4pt}
\setlength{\parindent}{0pt}
\EnableBpAbbreviations

\title{\textbf{Curso intensivo de los capítulos teóricos de la tesis}\\[4pt]
  \large Esquemas de demostración de los capítulos 7 a 10:\\
  cómo se interpretan y cómo se prueban}
\author{Material de apoyo para la defensa}
\date{}

\begin{document}
\maketitle
\thispagestyle{empty}

\begin{abstract}
\noindent
Este documento acompaña la lectura de los capítulos teóricos de la tesis:
el 7 (cálculo~\(\lambda\) y tipos simples), el 8 (cubo~\(\lambda\)), el 9
(deducción natural, semántica y Curry--Howard) y el 10 (Cálculo de
Construcciones Inductivas). Los capítulos 1 a 6 son de contexto y no
contienen demostraciones; los capítulos 11 y 12 tienen su propio curso.

Aquí no aparece Lean. Todo está escrito en matemática de papel, con la
notación de la tesis. El núcleo es un catálogo de \emph{esquemas de
demostración}: las formas de razonamiento que se repiten en esos capítulos.
Cada esquema tiene la misma estructura: su forma, por qué es válido, cómo se
redacta, dónde lo usa la tesis y qué hay que vigilar. La numeración (Lema~7.23,
Proposición~9.9, etc.) es la de la versión compilada de la tesis.
\end{abstract}

\tableofcontents
\newpage

%=====================================================================
\section{Cómo usar este documento}
%=====================================================================

Las demostraciones de los capítulos 7 a 10 usan un repertorio pequeño de
técnicas. Casi todas son variantes de tres ideas:

\begin{enumerate}
  \item \textbf{Inducción.} Los objetos de estos capítulos (términos, tipos,
    fórmulas, derivaciones, números naturales) se definen inductivamente, es
    decir, como lo mínimo cerrado bajo unas reglas. Por eso casi todo se
    demuestra recorriendo esas reglas una por una.
  \item \textbf{Construcción.} Probar un juicio \(\Gamma\vdash M:A\) o un
    secuente \(\Gamma\vdash\varphi\) es \emph{construir} una derivación. Bajo
    Curry--Howard, construir una prueba y escribir un programa son la misma
    actividad.
  \item \textbf{Reducción a un resultado anterior.} Varios corolarios (la
    unicidad de la forma normal, Church--Rosser para la conversión) se
    obtienen combinando un teorema con una propiedad sencilla de los objetos
    involucrados.
\end{enumerate}

Para cada resultado conviene distinguir tres situaciones, porque en una
defensa se pregunta por ellas:

\begin{center}
\begin{tabular}{@{}ll@{}}
\dem & la tesis da la demostración completa;\\[2pt]
\esb & la tesis da un esbozo y remite a la fuente para los detalles;\\[2pt]
\cit & la tesis enuncia el resultado y lo toma de la bibliografía.
\end{tabular}
\end{center}

\begin{nota}[Convención de lectura]
Los recuadros verdes son esquemas. Los grises reproducen o resumen un
argumento de la tesis. Los naranjas señalan los puntos delicados: pasos que la
tesis da por obvios y que un examinador puede pedir que se justifiquen. Cuando
este documento añade un argumento que la tesis no escribe (por ejemplo, el
contraejemplo al diamante de \(\bred\)), lo indica como
\emph{complemento}.
\end{nota}

%=====================================================================
\section{Mapa de los resultados}
%=====================================================================

La tabla reúne los resultados con demostración o esbozo de los capítulos 7
a 10, el esquema que usan (numerado como en la Sección~\ref{sec:catalogo}) y
su estado.

{\small
\begin{longtable}{@{}p{0.13\textwidth}p{0.44\textwidth}p{0.17\textwidth}p{0.16\textwidth}@{}}
\toprule
\textbf{Resultado} & \textbf{Enunciado (resumido)} & \textbf{Esquemas} & \textbf{Estado}\\
\midrule
\endhead
Lema 7.6 & \(M\in\Sub(M)\); transitividad de \(\Sub\) & 1, 2 & \dem\\
Lema 7.23 & Lema de sustitución & 1, 2, 3, 4 & \dem\\
Lema 7.31 & \(\bmred\) extiende \(\bred\), es reflexiva y transitiva & 5 & \dem\\
Lema 7.32 & \(\beq\) extiende \(\bmred\) y es de equivalencia & 5 & \dem\\
Teorema 7.37 & Church--Rosser (confluencia) & 8 & \esb\\
Lemas 7.40--7.41 & Sustitución y diamante para \(\pr\) & 1, 9 & \enu\\
Corolario 7.43 & Church--Rosser para la conversión & 6 & \dem\\
Corolario 7.44 & Unicidad de la forma normal & 7 & \dem\\
Ejemplo 7.49 & \(x\,x\) no es tipable & 17 & \dem\\
Ejemplos 7.52--7.53 & Derivaciones en \(\lambda\!\to\) & 10 & \dem\\
Lema 7.57 & \(\FV(L)\subseteq\dom(\Gamma)\) & 9 & \dem\\
Teorema 7.59 & Decidibilidad en \(\lambda\!\to\) & 23 & \cit\\
Teorema 7.63 & Church--Rosser en \(\lambda\!\to\) & 8, 23 & \cit\\
Ejemplo 8.5 & Identidad polimórfica en \(\lambda2\) & 10 & \dem\\
Ejemplos 8.13, 8.15 & Variable y debilitamiento en \(\lambda\underline{\omega}\) & 11 & \dem\\
Ejemplos 8.20, 8.22 & Derivación abreviada; conversión & 11 & \dem\\
§8.3.3 & Un habitante de \((\Pi x\,\Pi y.\,Q\,x\,y)\to\Pi u.\,Q\,u\,u\) & 12, 18 & \dem\\
Teorema 8.27 & Propiedades del cubo & 23 & \cit\\
Proposición 9.6 & \(p\wedge(q\vee r)\vdash(p\wedge q)\vee(p\wedge r)\) & 13 & \dem\\
Proposición 9.9 & \(p\to q\dashv\vdash\neg q\to\neg p\) & 13, 14 & \dem\\
Teorema 9.13 & Solidez & 9, 15 & \esb\\
Teorema 9.14 & Completitud (Kalmár) & 16 & \esb\\
Ejemplo 9.17 & Completitud como construcción: \(p\to p\) & 16 & \dem\\
Ejemplo 9.20 & Tautologías bajo BHK & 17, 18 & \dem\\
Teorema 9.22 & Isomorfismo de Curry--Howard & 18, 23 & \cit\\
Ejemplo 9.24 & La prueba como función \emph{swap} & 18 & \dem\\
Def.~9.25 & Normalización como evaluación & 19 & ---\\
§10.2.3 & Iteración sin inducción (Geuvers) & 21, 23 & \cit\\
§10.3.1--10.3.5 & Recursor, eliminación e \(\iota\)-reducción & 20 & ---\\
§10.3.2 & Positividad estricta & 22 & ---\\
\bottomrule
\end{longtable}
}

%=====================================================================
\section{Lo básico: definiciones inductivas y sus dos usos}
%=====================================================================

\subsection{Conjuntos definidos por reglas}

La Definición 7.1 introduce \(\Lambda\) con tres reglas: toda variable es un
término; si \(M,N\in\Lambda\) entonces \((MN)\in\Lambda\); si \(u\in V\) y
\(M\in\Lambda\) entonces \((\lambda u.M)\in\Lambda\). Lo implícito, y lo que
hace funcionar todas las pruebas, es la cláusula de \emph{minimalidad}:

\begin{center}
\(\Lambda\) es el \textbf{menor} conjunto que contiene a las variables y es cerrado bajo
aplicación y abstracción.
\end{center}

El mismo molde aparece en toda la parte teórica:

\begin{center}
\renewcommand{\arraystretch}{1.25}
\begin{tabular}{@{}lll@{}}
\toprule
\textbf{Objeto} & \textbf{Reglas} & \textbf{Dónde}\\
\midrule
términos \(\Lambda\), \(\Lambda_{\mathbb T}\) & variable, aplicación, abstracción & Def.~7.1, 7.50\\
tipos simples \(\mathbb T\) & variable de tipo, flecha & Def.~7.45\\
fórmulas \(\Phi\) & \(\bot\), átomo, \(\to\), \(\vee\), \(\wedge\) & Def.~9.21\\
\(\bred\) & base \((\lambda x.M)N\bred M[x:=N]\) y compatibilidad & Def.~7.24\\
\(\Gamma\vdash M:\sigma\) & (variable), (aplicación), (abstracción), \dots & §7.3.3, cap.~8\\
\(\Gamma\vdash\varphi\) & reglas de introducción y eliminación & §9.1, §9.4.2\\
\(\mathbb N\), listas, \dots & constructores (\(0\), \(S\)) & §10.3\\
\bottomrule
\end{tabular}
\end{center}

Una definición inductiva sirve para dos cosas, y cada una da un esquema de
demostración:

\begin{itemize}
  \item \textbf{Para construir:} un objeto pertenece al conjunto si se puede
    obtener aplicando las reglas un número finito de veces. Para probar
    \(\Gamma\vdash M:\sigma\) basta exhibir esa sucesión de reglas, es decir,
    una derivación.
  \item \textbf{Para razonar sobre todos los objetos:} si una propiedad se
    conserva bajo cada regla, la minimalidad garantiza que vale para todo el
    conjunto. Esto es la \emph{inducción}: estructural si el conjunto son
    términos, sobre derivaciones si son juicios.
\end{itemize}

El Capítulo~10 convierte esta observación en un teorema del sistema: al
declarar un tipo inductivo se generan sus constructores (las reglas) y su
eliminador (el principio de inducción); véase el Esquema~20.

\subsection{Notación que conviene tener presente}

\begin{itemize}
  \item \(M\equiv N\): identidad sintáctica. Desde la Sección~7.2.4 la tesis
    identifica términos \(\alpha\)-equivalentes. Por eso \(\equiv\) entre
    resultados de sustituciones se lee \emph{módulo \(\alpha\)}, y se
    escribe \(\aeq\) cuando interesa mostrar el renombramiento.
  \item \(\bred\): un paso de reducción \(\beta\). \(\bmred\): cero o más
    pasos. \(\beq\): conversión, es decir, cadenas que pueden ir en los dos
    sentidos.
  \item \(\Gamma\vdash M:\sigma\): juicio de tipado. \(\Gamma\vdash\varphi\):
    secuente de deducción natural. \(\Gamma\vDash\varphi\): consecuencia
    semántica. El mismo símbolo \(\vdash\) se usa en los dos primeros con
    significados distintos (la tesis lo advierte en §10.3.3).
\end{itemize}

%=====================================================================
\section{Catálogo de esquemas}\label{sec:catalogo}
%=====================================================================

\subsection*{Parte A. Inducción y análisis de casos sobre términos}

%---------------------------------------------------------------------
\begin{esquema}{Esquema 1. Inducción estructural sobre términos}\label{esq:estructural}
\forma Para probar \(\forall M\in\Lambda.\ \Phi(M)\) se prueban tres cosas:
\begin{enumerate}[label=(\roman*)]
  \item \textbf{caso variable:} \(\Phi(x)\) para toda variable \(x\);
  \item \textbf{caso aplicación:} si \(\Phi(P)\) y \(\Phi(Q)\) (hipótesis de
    inducción), entonces \(\Phi(PQ)\);
  \item \textbf{caso abstracción:} si \(\Phi(P)\), entonces
    \(\Phi(\lambda x.P)\).
\end{enumerate}
\porque El conjunto \(\{M\mid\Phi(M)\}\) contiene a las variables y es cerrado
bajo aplicación y abstracción. Como \(\Lambda\) es el menor conjunto con esas
propiedades, \(\Lambda\subseteq\{M\mid\Phi(M)\}\).

\plantilla «Se procede por inducción sobre la estructura de \(M\).
\emph{Caso \(M\equiv x\):} \dots\ \emph{Caso \(M\equiv PQ\):} por hipótesis de
inducción, \(\Phi(P)\) y \(\Phi(Q)\); entonces \dots\ \emph{Caso
\(M\equiv\lambda x.P\):} por hipótesis de inducción, \(\Phi(P)\); entonces
\dots»

\donde Lema 7.6 (transitividad de \(\Sub\)), Lema 7.23 (sustitución) y, en
general, cualquier propiedad de \(\Sub\), \(\FV\) o de la sustitución. La
Nota~7.5 de la tesis presenta el método.
\end{esquema}

\begin{ejemplotesis}{Lema 7.6 (transitividad): inducción sobre \(N\)}
\textbf{Enunciado.} Si \(L\in\Sub(M)\) y \(M\in\Sub(N)\), entonces
\(L\in\Sub(N)\).

\textbf{Elección de la variable.} Se fijan \(L\) y \(M\) y se hace inducción
sobre \(N\) con la propiedad
\[
\Phi(N):\quad L\in\Sub(M)\ \wedge\ M\in\Sub(N)\ \Longrightarrow\ L\in\Sub(N).
\]
\textbf{Caso \(N\equiv PQ\).} Por definición,
\(\Sub(PQ)=\Sub(P)\cup\Sub(Q)\cup\{PQ\}\), así que \(M\in\Sub(N)\) deja tres
posibilidades:
\begin{itemize}
  \item \(M\in\Sub(P)\): la hipótesis de inducción \(\Phi(P)\) da
    \(L\in\Sub(P)\subseteq\Sub(N)\);
  \item \(M\in\Sub(Q)\): simétrico, con \(\Phi(Q)\);
  \item \(M\equiv PQ\equiv N\): entonces \(L\in\Sub(M)=\Sub(N)\) directamente.
\end{itemize}
Los casos variable y abstracción son análogos (uno y dos subcasos).
\end{ejemplotesis}

\begin{trampa}
\begin{itemize}
  \item \textbf{¿Por qué sobre \(N\) y no sobre \(M\)?} La hipótesis que hay
    que descomponer es \(M\in\Sub(N)\), y \(\Sub(N)\) se calcula según la
    forma de \(N\). Conocer la forma de \(M\) no dice nada sobre \(\Sub(N)\).
  \item \textbf{La reflexividad no usa la hipótesis de inducción.} Cada
    cláusula de \(\Sub\) añade el propio término (\(\{x\}\), \(\cup\{MN\}\),
    \(\cup\{\lambda x.M\}\)), así que \(M\in\Sub(M)\) sale directamente de la
    definición en cada caso. La tesis la presenta caso por caso dentro del
    marco inductivo, pero basta un análisis de casos (Esquema~2).
  \item Aquí \(L\) y \(M\) se mantienen fijos, y la hipótesis de inducción se
    aplica con los mismos \(L,M\). En otras pruebas hace falta
    \emph{generalizar}, es decir, cuantificar una variable antes de inducir
    (véase el Esquema~9 con el contexto \(\Gamma\)).
\end{itemize}
\end{trampa}

%---------------------------------------------------------------------
\begin{esquema}{Esquema 2. Análisis de casos según la definición}
\forma Cuando un objeto se define por cláusulas (como \(\Sub\), \(\FV\) o la
sustitución), una pertenencia o una igualdad se descompone en tantos casos
como cláusulas apliquen, y cada caso se cierra con la cláusula
correspondiente.

\porque Las cláusulas son exhaustivas (todo término es una variable, una
aplicación o una abstracción) y mutuamente excluyentes. La disyunción de los
casos es verdadera y cada rama prueba la meta: es la eliminación de la
disyunción (\(\vee\)e del Capítulo~9) aplicada a la matemática informal.

\plantilla «Se distinguen tres subcasos según la variable \(z\): \(z\equiv x\),
\(z\equiv y\) y \(z\not\equiv x,y\). En el primero \dots»

\donde El caso base del Lema 7.23. La pertenencia
\(M\in\Sub(PQ)\iff M\in\Sub(P)\vee M\in\Sub(Q)\vee M\equiv PQ\) del Lema 7.6.
\end{esquema}

\begin{ejemplotesis}{Lema 7.23, caso base: tres subcasos}
Con \(x\not\equiv y\) y \(x\notin\FV(L)\) se quiere
\(z[x:=N][y:=L]\equiv z[y:=L][x:=N[y:=L]]\) para una variable \(z\).
\begin{itemize}
  \item \(z\equiv x\): ambos lados valen \(N[y:=L]\). El lado derecho usa
    \(x[y:=L]\equiv x\), que requiere \(x\not\equiv y\).
  \item \(z\equiv y\): el lado izquierdo es \(y[y:=L]\equiv L\). El derecho es
    \(L[x:=N[y:=L]]\equiv L\), que requiere \(x\notin\FV(L)\).
  \item \(z\) distinta de ambas: los dos lados valen \(z\).
\end{itemize}
Cada hipótesis del lema se usa en exactamente un subcaso. Es una buena forma
de explicar \emph{por qué} el lema tiene esas hipótesis.
\end{ejemplotesis}

\begin{trampa}[Un lema auxiliar implícito]
El subcaso \(z\equiv y\) usa el hecho
\[
x\notin\FV(L)\ \Longrightarrow\ L[x:=K]\equiv L,
\]
que la tesis no enuncia por separado. Se demuestra por inducción estructural
sobre \(L\) (Esquema~1). Si un examinador pregunta por esa línea, esta es la
respuesta.
\end{trampa}

%---------------------------------------------------------------------
\begin{esquema}{Esquema 3. Cadena ecuacional justificada}
\forma Una igualdad (o identidad \(\equiv\)) se prueba con una cadena
\(A\equiv A_1\equiv\cdots\equiv B\) en la que cada eslabón lleva una
justificación: una cláusula de una definición, la hipótesis de inducción o un
lema.

\porque \(\equiv\) (y \(\aeq\), \(\beq\)) es reflexiva, simétrica y
transitiva, así que la cadena prueba \(A\equiv B\).

\plantilla Una línea por paso, con la razón a la derecha:
\[
\begin{aligned}
(PQ)[x:=N][y:=L]
&\equiv (P[x:=N][y:=L])\,(Q[x:=N][y:=L])\\
&\qquad\text{(def. de sustitución, dos veces)}\\
&\equiv (P[y:=L][x:=N[y:=L]])\,(Q[y:=L][x:=N[y:=L]])\\
&\qquad\text{(hipótesis de inducción en \(P\) y en \(Q\))}\\
&\equiv (PQ)[y:=L][x:=N[y:=L]]\\
&\qquad\text{(def. de sustitución, leída al revés)}
\end{aligned}
\]

\donde Caso aplicación del Lema 7.23. Ejemplo 7.22 (cálculo de sustituciones
regla por regla). Reducciones en los Ejemplos 7.26, 7.30 y 7.42.
\end{esquema}

\begin{nota}[El paso «leído al revés»]
El último eslabón usa la definición de sustitución de derecha a izquierda:
se reconoce que \((P'Q')[x:=K]\) \emph{se escribe} \(P'[x:=K]\,Q'[x:=K]\).
Es el momento en que la prueba «vuelve a empaquetar» el resultado, y conviene
señalarlo al explicarla.
\end{nota}

%---------------------------------------------------------------------
\begin{esquema}{Esquema 4. Variables ligadas: elegir un representante fresco}
\forma En el caso de abstracción, antes de calcular se renombra la variable
ligada a una \(z\) \emph{fresca}: distinta de todas las variables relevantes y
fuera de las variables libres de los términos que se sustituyen.

\porque La tesis identifica términos \(\alpha\)-equivalentes (§7.2.4, tras la
Definición 7.19) y adopta la convención de Barendregt (§7.2.6). Renombrar una
variable ligada no cambia el término módulo \(\alpha\), y elegir \(z\) fresca
hace que la cláusula (3) de la sustitución pase \emph{sin} renombrar, es decir,
\((\lambda z.P)[x:=N]\equiv\lambda z.(P[x:=N])\).

\plantilla «Se puede renombrar \(z\) para que \(z\not\equiv x\),
\(z\not\equiv y\) y \(z\notin\FV(N)\cup\FV(L)\). Entonces la sustitución entra
bajo la abstracción y \dots»

\donde Caso abstracción del Lema 7.23. Ejemplo 7.22, donde no renombrar
produce la captura de la variable \(y\).
\end{esquema}

\begin{trampa}[Qué hay que comprobar en el caso abstracción del Lema 7.23]
La tesis escribe el caso en tres líneas. Cada línea usa una de las condiciones
de frescura:
\begin{enumerate}
  \item \((\lambda z.P)[x:=N]\equiv\lambda z.P[x:=N]\) usa \(z\not\equiv x\) y
    \(z\notin\FV(N)\);
  \item aplicar después \([y:=L]\) usa \(z\not\equiv y\) y \(z\notin\FV(L)\);
  \item en el lado derecho, aplicar \([x:=N[y:=L]]\) usa
    \(z\notin\FV(N[y:=L])\). Esto se sigue de
    \(\FV(N[y:=L])\subseteq\FV(N)\cup\FV(L)\), otra pequeña inducción
    estructural.
\end{enumerate}
Como la tesis trabaja módulo \(\alpha\), el Lema 7.23 afirma una identidad
módulo \(\alpha\). Una igualdad sintáctica literal no se cumple en general,
porque la elección de la variable fresca no es única.
\end{trampa}

\subsection*{Parte B. Relaciones de reducción}

%---------------------------------------------------------------------
\begin{esquema}{Esquema 5. Relaciones definidas por cadenas finitas}
\forma \(M\bmred N\) y \(M\beq N\) se definen por la \emph{existencia} de una
cadena finita \(M\equiv M_0,\dots,M_n\equiv N\) (Definiciones 7.27 y 7.28).
\begin{itemize}
  \item Para \emph{probar} una de estas relaciones se exhibe una cadena
    (introducción del \(\exists\)).
  \item Para \emph{usarla} se toma una cadena cualquiera que la testifique
    (eliminación del \(\exists\)) y se razona sobre ella.
\end{itemize}

\porque La definición es un existencial: las dos operaciones son exactamente
sus reglas de introducción y eliminación.

\plantilla Las cuatro operaciones del Lema 7.31 y del Lema 7.32:
\begin{center}
\begin{tabular}{@{}ll@{}}
\toprule
\textbf{Propiedad} & \textbf{Cadena que se exhibe}\\
\midrule
extensión (\(\bred\subseteq\bmred\)) & la cadena de un paso (\(k=1\))\\
reflexividad & la cadena vacía (\(k=0\))\\
transitividad & la concatenación (\(r+s\) pasos)\\
simetría de \(\beq\) & la cadena invertida\\
\bottomrule
\end{tabular}
\end{center}

\donde Lema 7.31 y Lema 7.32.
\end{esquema}

\begin{nota}[Por qué la simetría solo vale para \(\beq\)]
La inversión funciona porque cada eslabón de \(\beq\) es
\(M_i\bred M_{i+1}\) \emph{o} \(M_{i+1}\bred M_i\). Al invertir la cadena, un
eslabón hacia adelante se convierte en uno hacia atrás, y viceversa. En
\(\bmred\) todos los eslabones van hacia adelante, así que invertir la cadena
no produce otra cadena de \(\bmred\). Por ejemplo,
\((\lambda x.x)y\bmred y\) pero no \(y\bmred(\lambda x.x)y\).
\end{nota}

\begin{nota}[Una presentación equivalente]
\(\bmred\) es la menor relación reflexiva y transitiva que contiene a
\(\bred\), y \(\beq\) es la menor relación de equivalencia que contiene a
\(\bred\). Con esa presentación, el Lema 7.31 y la segunda parte del
Lema 7.32 se cumplen por construcción. La tesis elige la presentación con
cadenas porque hace visibles los testigos.
\end{nota}

%---------------------------------------------------------------------
\begin{esquema}{Esquema 6. Inducción sobre la longitud de una cadena}
\forma Para probar algo de toda cadena \(M_0\bconv M_1\bconv\cdots\bconv M_k\)
se hace inducción sobre \(k\in\mathbb N\): el caso \(k=0\) y, para el paso
inductivo, se aplica la hipótesis a los primeros \(k-1\) eslabones y se trata
el último por casos.

\porque Es inducción sobre los naturales aplicada a la longitud de la cadena.

\donde Corolario 7.43 (Church--Rosser para la conversión).
\end{esquema}

\begin{ejemplotesis}{Corolario 7.43: del diamante para \(\bmred\) a la conversión}
\textbf{Meta.} Si \(M\beq N\), existe \(L\) con \(M\bmred L\) y \(N\bmred L\).

\textbf{\(k=0\).} \(M\equiv N\); se toma \(L\equiv M\).

\textbf{Paso.} La hipótesis de inducción da \(L'\) con \(M\bmred L'\) y
\(M_{k-1}\bmred L'\). El último eslabón tiene dos orientaciones:
\begin{itemize}
  \item \(M_k\bred M_{k-1}\): entonces \(M_k\bmred L'\) por transitividad
    (Lema 7.31), y \(L\equiv L'\).
  \item \(M_{k-1}\bred M_k\): \(M_{k-1}\) se reduce a \(L'\) y a \(M_k\).
    Church--Rosser (Teorema 7.37) da \(L\) con \(L'\bmred L\) y
    \(M_k\bmred L\); además \(M\bmred L'\bmred L\).
\end{itemize}
\[
\begin{tikzcd}[column sep=small, row sep=small]
M \arrow[rd, two heads] & & M_{k-1} \arrow[ld, two heads] \arrow[rd] & \\
& L' \arrow[rd, two heads, dashed] & & M_k \arrow[ld, two heads, dashed]\\
& & L &
\end{tikzcd}
\]
\end{ejemplotesis}

%---------------------------------------------------------------------
\begin{esquema}{Esquema 7. Unicidad vía confluencia y objetos extremales}
\forma Para probar que un objeto con cierta propiedad es único: se suponen dos
(\(N_1,N_2\)), se usa un teorema que los conecta (confluencia) y se usa que
los objetos con esa propiedad no pueden «moverse» (una forma normal solo se
reduce a sí misma).

\porque Si \(N_1\bmred L\) y \(N_1\) no tiene redexes, la única cadena posible
desde \(N_1\) tiene longitud \(0\), así que \(N_1\equiv L\). Lo mismo vale
para \(N_2\), y por tanto \(N_1\equiv L\equiv N_2\).

\plantilla «Sean \(N_1,N_2\) dos formas normales de \(M\). Por simetría y
transitividad \(N_1\beq N_2\); por el Corolario 7.43 existe \(L\) con
\(N_1\bmred L\) y \(N_2\bmred L\). Como las formas normales no tienen redexes,
\(N_1\equiv L\equiv N_2\).»

\donde Corolario 7.44 (unicidad de la forma normal).
\end{esquema}

\begin{trampa}
\begin{itemize}
  \item La forma normal se define con \(M\beq N\) (Definición 7.33), no con
    \(M\bmred N\). Por eso la tesis necesita primero el Corolario 7.43, que
    extiende la confluencia de \(\bmred\) a \(\beq\). Es la razón de ser de
    ese corolario.
  \item Church--Rosser \textbf{no} dice que todo término tenga forma normal:
    \(\Omega\) no la tiene (Ejemplo 7.36). Solo dice que, si existe, es
    única. La existencia, para términos tipados, es la normalización fuerte
    (Teorema 8.27, citado).
\end{itemize}
\end{trampa}

%---------------------------------------------------------------------
\begin{esquema}{Esquema 8. Confluencia por el método de Tait y Martin-Löf}\label{esq:tml}
\forma Para probar que \(\bmred\) tiene la propiedad del diamante:
\begin{enumerate}
  \item se define una relación auxiliar \(\pr\) (reducción paralela,
    Definición 7.39) que contrae cualquier conjunto de redexes a la vez;
  \item se prueba que \(\pr\) tiene la propiedad del diamante \emph{en un
    paso} (Lema 7.41), con ayuda del lema de sustitución para \(\pr\)
    (Lema 7.40);
  \item se verifica \(\bred\ \subseteq\ \pr\ \subseteq\ \bmred\). Como
    consecuencia, la clausura reflexiva y transitiva de \(\pr\) coincide con
    \(\bmred\);
  \item el diamante de \(\pr\) pasa a su clausura transitiva «embaldosando»
    el rectángulo.
\end{enumerate}

\porque El paso 4 es puramente combinatorio. Si \(M\pr^{*}N_1\) en \(m\) pasos
y \(M\pr^{*}N_2\) en \(n\) pasos, se completa una cuadrícula de
\(m\times n\) diamantes de un paso:
\begin{center}
\begin{tikzpicture}[x=1.45cm,y=1.0cm,>=Stealth]
  \foreach \i in {0,1,2} \foreach \j in {0,1} {
    \node (n\i\j) at (\i,-\j) {\(\bullet\)};
  }
  \node[above left] at (0,0) {\(M\)};
  \node[above right] at (2,0) {\(N_1\)};
  \node[below left] at (0,-1) {\(N_2\)};
  \node[below right] at (2,-1) {\(L\)};
  \foreach \i/\k in {0/1,1/2} {
    \draw[->] (n\i0) -- node[above]{\scriptsize\(\pr\)} (n\k0);
    \draw[->,dashed] (n\i1) -- node[below]{\scriptsize\(\pr\)} (n\k1);
  }
  \foreach \i in {0} \draw[->] (n\i0) -- node[left]{\scriptsize\(\pr\)} (n\i1);
  \foreach \i in {1,2} \draw[->,dashed] (n\i0) -- node[right]{\scriptsize\(\pr\)} (n\i1);
\end{tikzpicture}
\end{center}
Cada cuadrado se cierra con el Lema 7.41. Las flechas discontinuas existen
por el diamante, y la esquina \(L\) es el reducto común.

\donde Teorema 7.37 (esbozo, siguiendo a Barendregt, Teorema 3.2.8) y
Teorema 7.63 (mismo argumento en \(\lambda\!\to\)).
\end{esquema}

\begin{trampa}[¿Por qué no se embaldosa directamente con \(\bred\)? (complemento)]
Porque \(\bred\) \emph{no} tiene el diamante en un paso. Sea
\(M\equiv(\lambda x.xx)\,((\lambda y.y)z)\):
\[
M\bred((\lambda y.y)z)\,((\lambda y.y)z),\qquad M\bred(\lambda x.xx)\,z.
\]
El segundo llega a \(zz\) en un paso, pero el primero necesita dos, porque el
redex interno quedó \emph{duplicado}. La reducción paralela contrae las dos
copias a la vez, y por eso con \(\pr\) los cuadrados sí se cierran en un paso.
Este es el motivo técnico de introducir \(\pr\).

Tampoco basta la confluencia \emph{local} (cerrar cada par de pasos con
\(\bmred\)). Sin terminación no implica confluencia (lema de Newman), y el
cálculo no tipado no termina.
\end{trampa}

\subsection*{Parte C. Juicios y derivaciones}

%---------------------------------------------------------------------
\begin{esquema}{Esquema 9. Inducción sobre derivaciones (inducción por reglas)}\label{esq:derivaciones}
\forma Para probar «si \(\Gamma\vdash L:\sigma\) es derivable, entonces
\(\Psi(\Gamma,L,\sigma)\)», se recorre cada regla de tipado: suponiendo
\(\Psi\) para las premisas, se prueba \(\Psi\) para la conclusión.

\porque La derivabilidad es la menor relación cerrada bajo las reglas. Si
\(\Psi\) también es cerrada bajo ellas, contiene a la derivabilidad. Es el
Esquema~1 aplicado a \emph{árboles de derivación} en lugar de términos.

\plantilla «Por inducción sobre la derivación de \(\Gamma\vdash L:\sigma\).
\emph{Regla (variable):} \dots\ \emph{Regla (aplicación):} por hipótesis
inductiva aplicada a las dos premisas \dots\ \emph{Regla (abstracción):} por
hipótesis inductiva aplicada a la premisa \(\Gamma,x:\sigma\vdash M:\tau\)
\dots»

\donde Lema 7.57 (variables libres). Teorema 9.13 (solidez), donde Huth y Ryan
lo llaman «inducción sobre la longitud de la prueba».
\end{esquema}

\begin{ejemplotesis}{Lema 7.57: \(\FV(L)\subseteq\dom(\Gamma)\)}
\begin{itemize}
  \item \textbf{(variable)} \(x:\sigma\in\Gamma\) da \(x\in\dom(\Gamma)\) y
    \(\FV(x)=\{x\}\).
  \item \textbf{(aplicación)} Por las hipótesis inductivas,
    \(\FV(MN)=\FV(M)\cup\FV(N)\subseteq\dom(\Gamma)\).
  \item \textbf{(abstracción)} La premisa vive en el contexto
    \emph{extendido} \(\Gamma,x:\sigma\). La hipótesis inductiva da
    \(\FV(M)\subseteq\dom(\Gamma)\cup\{x\}\), y entonces
    \(\FV(\lambda x{:}\sigma.M)=\FV(M)\setminus\{x\}\subseteq\dom(\Gamma)\).
\end{itemize}
\end{ejemplotesis}

\begin{trampa}
\begin{itemize}
  \item \textbf{No fijar el contexto.} En la regla (abstracción) la premisa
    tiene un contexto distinto del de la conclusión. La propiedad \(\Psi\)
    debe enunciarse para \emph{todo} \(\Gamma\). Si se fija \(\Gamma\) antes
    de inducir, la hipótesis inductiva no se aplica a la premisa.
  \item \textbf{Inducción estructural frente a inducción sobre derivaciones.}
    En la primera los casos son las formas del término; en la segunda, las
    reglas. En \(\lambda\!\to\) coinciden porque el sistema está dirigido por
    la sintaxis: cada forma de término admite una sola regla. En
    \(\lambda\underline{\omega}\), con (debilitamiento) y (conversión), ya no
    coinciden: el mismo término puede obtenerse por varias reglas.
\end{itemize}
\end{trampa}

%---------------------------------------------------------------------
\begin{esquema}{Esquema 10. Construir una derivación de tipado}
\forma Se trabaja \emph{de la meta hacia las hipótesis}: la forma del término
indica la regla.
\begin{center}
\begin{tabular}{@{}lll@{}}
\toprule
\textbf{El término es\dots} & \textbf{Regla} & \textbf{Nueva meta}\\
\midrule
una variable \(x\) & (variable) & comprobar \(x:\sigma\in\Gamma\)\\
\(\lambda x{:}\sigma.M\) con tipo \(\sigma\to\tau\) & (abstracción) & \(\Gamma,x:\sigma\vdash M:\tau\)\\
\(MN\) con tipo \(\tau\) & (aplicación) & \(\Gamma\vdash M:\sigma\to\tau\) y \(\Gamma\vdash N:\sigma\)\\
\(\Lambda\alpha{:}*.M\) con tipo \(\Pi\alpha{:}*.A\) & (abstracción\(_2\)) & \(\Gamma,\alpha:*\vdash M:A\)\\
\bottomrule
\end{tabular}
\end{center}
Una vez llegada a las hojas, la derivación se escribe al revés: de las hojas a
la raíz.

\porque Cada regla se lee en los dos sentidos. Hacia abajo construye; hacia
arriba dice qué falta. Cuando se llega a hojas cerradas por (variable), el
árbol queda completo.

\donde Ejemplos 7.52 (función constante), 7.53 (\(\eta\)-expansión de la
identidad) y 8.5 (identidad polimórfica).
\end{esquema}

\begin{ejemplotesis}{Ejemplo 7.53 en los tres formatos de la tesis}
\textbf{Árbol.}
\begin{prooftree}
  \AxiomC{\((y:\alpha\to\beta)\in\Gamma_1\)}
  \RightLabel{\scriptsize(var)}
  \UnaryInfC{\(\Gamma_1\vdash y:\alpha\to\beta\)}
  \AxiomC{\((z:\alpha)\in\Gamma_1\)}
  \RightLabel{\scriptsize(var)}
  \UnaryInfC{\(\Gamma_1\vdash z:\alpha\)}
  \RightLabel{\scriptsize(apl)}
  \BinaryInfC{\(\Gamma_1\vdash y\,z:\beta\)}
  \RightLabel{\scriptsize(abs)}
  \UnaryInfC{\(y:\alpha\to\beta\vdash\lambda z{:}\alpha.\,y\,z:\alpha\to\beta\)}
  \RightLabel{\scriptsize(abs)}
  \UnaryInfC{\(\vdash\lambda y{:}\alpha\to\beta.\,\lambda z{:}\alpha.\,y\,z:(\alpha\to\beta)\to\alpha\to\beta\)}
\end{prooftree}
con \(\Gamma_1=\{y:\alpha\to\beta,\ z:\alpha\}\).

\textbf{Lineal.} Los juicios (1)--(5) numerados, cada uno con su regla y sus
premisas (es la redacción de la demostración en la tesis).

\textbf{Banderas.} Cada declaración abre una bandera; las líneas dentro de ella
la tienen en su contexto. Cerrar la bandera es aplicar (abstracción):
\begin{center}
\begin{flagderiv}[fdsiete]
  \introduce{y}{y:\alpha\to\beta}{}
  \introduce{z}{z:\alpha}{}
  \step{yy}{y:\alpha\to\beta}{(var)}
  \step{zz}{z:\alpha}{(var)}
  \step{yz}{y\,z:\beta}{(apl) en \ref{yy}, \ref{zz}}
  \conclude{lz}{\lambda z{:}\alpha.\,y\,z:\alpha\to\beta}{(abs) en \ref{yz}}
  \conclude{ly}{\lambda y{:}\alpha\to\beta.\,\lambda z{:}\alpha.\,y\,z:(\alpha\to\beta)\to\alpha\to\beta}{(abs) en \ref{lz}}
\end{flagderiv}
\end{center}
Los tres formatos codifican la misma información. El de banderas evita
repetir el contexto en cada línea, y su «caja» es la misma que la de la
deducción natural (Nota 9.5).
\end{ejemplotesis}

%---------------------------------------------------------------------
\begin{esquema}{Esquema 11. Derivar con contextos bien formados (sistemas con sorts)}
\forma En \(\lambda\underline{\omega}\) y en los sistemas del cubo, la buena
formación del contexto es parte de la prueba. Cada declaración \(x:A\) que se
añade exige derivar antes \(\Gamma\vdash A:s\). Las reglas que hacen ese
trabajo son:
\[
\text{(sort)}\ \ \varnothing\vdash *:\square\qquad
\frac{\Gamma\vdash A:s}{\Gamma,x:A\vdash x:A}\ \text{(var)}\qquad
\frac{\Gamma\vdash A:B\quad\Gamma\vdash C:s}{\Gamma,x:C\vdash A:B}\ \text{(deb)}
\]
y, para cambiar un tipo por otro \(\beta\)-equivalente,
\[
\frac{\Gamma\vdash A:B\quad\Gamma\vdash B':s}{\Gamma\vdash A:B'}\ \ (B\beq B')\ \text{(conversión)}.
\]

\porque Sin la segunda premisa del debilitamiento (\(\Gamma\vdash C:s\)) se
podrían añadir declaraciones sin sentido, como \(x:3\). Sin la segunda
premisa de la conversión se podría «convertir» a una expresión mal formada
(la Definición 8.21 lo advierte).

\plantilla Formato de banderas, indicando en cada paso \emph{las dos}
premisas: «(debilitamiento) de 2, 3». La tesis señala que pueden coincidir
(«de 1, 1»), porque cada premisa cumple un papel distinto.

\donde Ejemplo 8.13 (\(\alpha:*,x:\alpha\vdash x:\alpha\)), Ejemplo 8.15
(debilitamiento), Ejemplo 8.20 (derivación completa frente a abreviada) y
Ejemplo 8.22 (conversión).
\end{esquema}

\begin{ejemplotesis}{Ejemplo 8.22: usar la conversión}
Con \(\Gamma\equiv\beta:*,\ x:(\lambda\alpha{:}*.\,\alpha\to\alpha)\,\beta\) se
quiere \(\Gamma\vdash x:\beta\to\beta\).
\begin{enumerate}
  \item \(\Gamma\vdash x:(\lambda\alpha{:}*.\,\alpha\to\alpha)\,\beta\) por
    (variable).
  \item \(\Gamma\vdash\beta\to\beta:*\) por (formación), a partir de
    \(\Gamma\vdash\beta:*\).
  \item \((\lambda\alpha{:}*.\,\alpha\to\alpha)\,\beta\bred\beta\to\beta\): un
    paso \(\beta\), ahora \emph{en el nivel de los tipos}.
  \item (conversión) con 1, 2 y 3.
\end{enumerate}
El punto conceptual es que el tipado se hace \emph{módulo cómputo}: dos tipos
que se reducen a lo mismo son intercambiables.
\end{ejemplotesis}

\begin{nota}[Derivaciones abreviadas (Ejemplo 8.20)]
Una derivación completa en \(\lambda\underline{\omega}\) tiene muchos pasos
rutinarios de (sort), (variable) y (debilitamiento). La versión abreviada los
omite y conserva solo (formación), (abstracción), (aplicación) y
(conversión). Es la misma convención que en matemática ordinaria omite
«\(x\) es un número real» en cada línea. Al explicar una derivación abreviada
conviene decir qué pasos se omitieron y por qué son rutinarios.
\end{nota}

%---------------------------------------------------------------------
\begin{esquema}{Esquema 12. Habitar un tipo es construir una prueba}
\forma Para probar una proposición \(A\) se busca un término \(M\) con
\(\Gamma\vdash M:A\) (construcción de términos, §7.3.4). Cada conectivo indica
qué hacer:
\begin{center}
\begin{tabular}{@{}lll@{}}
\toprule
\textbf{Meta / hipótesis} & \textbf{Qué se hace} & \textbf{Regla}\\
\midrule
meta \(A\to B\) & suponer \(z:A\) y habitar \(B\) & (abstracción)\\
meta \(\Pi u{:}S.\,C\) & tomar \(u:S\) arbitrario y habitar \(C\) & (abstracción)\\
hipótesis \(z:\Pi x{:}S.\,D\) & aplicar \(z\) a un \(a:S\) y obtener \(D[x:=a]\) & (aplicación)\\
hipótesis \(f:A\to B\) y \(a:A\) & aplicar y obtener \(f\,a:B\) & (aplicación)\\
\bottomrule
\end{tabular}
\end{center}

\porque Es la correspondencia PAT (Definiciones 7.61 y 8.25): la abstracción
es la introducción de \(\to\) y de \(\forall\) (generalización), y la
aplicación es su eliminación (modus ponens y especialización).

\donde §8.3.3 (lógica de predicados en \(\lambda P\)); Ejemplos 7.52 y 9.24.
\end{esquema}

\begin{ejemplotesis}{§8.3.3: si \(Q\) vale para todos los pares, vale para los pares \((u,u)\)}
Meta: habitar \((\Pi x{:}S.\,\Pi y{:}S.\,Q\,x\,y)\to(\Pi u{:}S.\,Q\,u\,u)\) con
\(S:*\) y \(Q:S\to S\to *\).
\begin{center}
\begin{flagderiv}[fdpi]
  \introduce{z}{z:\Pi x{:}S.\,\Pi y{:}S.\,Q\,x\,y}{supuesto}
  \introduce{u}{u:S}{$u$ arbitrario}
  \step{zu}{z\,u:\Pi y{:}S.\,Q\,u\,y}{(apl) en \ref{z}, \ref{u}}
  \step{zuu}{z\,u\,u:Q\,u\,u}{(apl) en \ref{zu}, \ref{u}}
  \conclude{lu}{\lambda u{:}S.\,z\,u\,u:\Pi u{:}S.\,Q\,u\,u}{(abs) en \ref{zuu}}
  \conclude{lz}{\lambda z.\,\lambda u{:}S.\,z\,u\,u:(\Pi x\,\Pi y.\,Q\,x\,y)\to\Pi u.\,Q\,u\,u}{(abs) en \ref{lu}}
\end{flagderiv}
\end{center}
\textbf{Lectura en español.} «Supóngase que \(Q(x,y)\) vale para todo par
(\(z\)). Sea \(u\in S\) arbitrario. Instanciando la hipótesis en \(x:=u\) e
\(y:=u\) se obtiene \(Q(u,u)\). Como \(u\) era arbitrario, \(\forall u.\,Q(u,u)\).»
El término \(\lambda z.\lambda u.\,z\,u\,u\) \emph{es} esa frase.
\end{ejemplotesis}

\subsection*{Parte D. Lógica: deducción natural y semántica}

%---------------------------------------------------------------------
\begin{esquema}{Esquema 13. Derivación en deducción natural con cajas}
\forma Una prueba es una lista numerada de fórmulas, cada una con su
justificación (premisa, suposición o regla aplicada a líneas anteriores). Las
reglas que \emph{descargan} una suposición (\(\to\)i, \(\neg\)i, \(\vee\)e y
RAA) abren una caja: empieza con la suposición y se cierra al aplicar la regla.

\porque Es la presentación de Fitch (o Huth--Ryan) de las reglas de §9.1. La
disciplina de cajas (Nota 9.5) garantiza que la conclusión final dependa solo
de las premisas.

\plantilla Una heurística que funciona casi siempre:
\begin{enumerate}
  \item mirar la \emph{meta}: si es \(A\to B\), abrir una caja con \(A\) y
    buscar \(B\) (\(\to\)i); si es \(\neg A\), abrir una caja con \(A\) y
    buscar \(\bot\) (\(\neg\)i);
  \item mirar las \emph{premisas}: descomponer las conjunciones (\(\wedge\)e) y
    usar las disyunciones por casos (\(\vee\)e);
  \item si nada de lo anterior funciona y la lógica es clásica, abrir una caja
    con \(\neg A\) y buscar \(\bot\) (RAA).
\end{enumerate}

\donde Proposiciones 9.6 y 9.9.
\end{esquema}

\begin{ejemplotesis}{Proposición 9.6: \(p\wedge(q\vee r)\vdash(p\wedge q)\vee(p\wedge r)\)}
\begin{center}
\begin{tabular}{lll}
1 & \(p\wedge(q\vee r)\) & premisa\\
2 & \(p\) & \(\wedge\)e\(_1\) 1\\
3 & \(q\vee r\) & \(\wedge\)e\(_2\) 1\\ \cline{2-3}
4 & \multicolumn{1}{|l}{\(q\)} & suposición (caso 1)\\
5 & \multicolumn{1}{|l}{\(p\wedge q\)} & \(\wedge\)i 2, 4\\
6 & \multicolumn{1}{|l}{\((p\wedge q)\vee(p\wedge r)\)} & \(\vee\)i\(_1\) 5\\ \cline{2-3}
7 & \multicolumn{1}{|l}{\(r\)} & suposición (caso 2)\\
8 & \multicolumn{1}{|l}{\(p\wedge r\)} & \(\wedge\)i 2, 7\\
9 & \multicolumn{1}{|l}{\((p\wedge q)\vee(p\wedge r)\)} & \(\vee\)i\(_2\) 8\\ \cline{2-3}
10 & \((p\wedge q)\vee(p\wedge r)\) & \(\vee\)e 3, 4--6, 7--9
\end{tabular}
\end{center}
Hay dos lecciones. La línea 2, fuera de las cajas, se usa dentro de ambas. Las
dos ramas tienen que terminar en \emph{la misma} fórmula; por eso las líneas 6
y 9 introducen la disyunción completa antes de cerrar.
\end{ejemplotesis}

\begin{trampa}
\begin{itemize}
  \item Lo que está dentro de una caja cerrada \emph{no} puede usarse
    después (líneas 4--9 no están disponibles en la línea 10, salvo como
    bloques citados por la regla).
  \item Bajo Curry--Howard esta restricción es el \emph{alcance de una
    variable ligada}. La suposición \(q\) de la línea 4 es la variable \(x\)
    de \(\operatorname{case}(L;x.M;y.N)\) (§9.5.1).
\end{itemize}
\end{trampa}

%---------------------------------------------------------------------
\begin{esquema}{Esquema 14. Equivalencia: dos direcciones y dónde entra lo clásico}
\forma \(\varphi\dashv\vdash\psi\) se prueba con dos derivaciones
independientes, \(\varphi\vdash\psi\) y \(\psi\vdash\varphi\)
(Definición 9.8). Conviene identificar en cuál de ellas, si en alguna, se usa
una regla clásica (RAA, LEM o \(\neg\neg\)e).

\porque Es la definición. Separar las direcciones permite ver cuál es
intuicionista y cuál no.

\donde Proposición 9.9 (contrapositiva).
\end{esquema}

\begin{ejemplotesis}{Proposición 9.9: \(p\to q\dashv\vdash\neg q\to\neg p\)}
\begin{minipage}[t]{0.49\textwidth}
\textbf{(\(\Rightarrow\)) intuicionista}\\[2pt]
\begin{tabular}{lll}
1 & \(p\to q\) & premisa\\ \cline{2-3}
2 & \multicolumn{1}{|l}{\(\neg q\)} & sup.\ (\(\to\)i)\\
3 & \multicolumn{1}{|l}{\quad\(p\)} & sup.\ (\(\neg\)i)\\
4 & \multicolumn{1}{|l}{\quad\(q\)} & \(\to\)e 1, 3\\
5 & \multicolumn{1}{|l}{\quad\(\bot\)} & \(\neg\)e 4, 2\\ \cline{3-3}
6 & \multicolumn{1}{|l}{\(\neg p\)} & \(\neg\)i 3--5\\ \cline{2-3}
7 & \(\neg q\to\neg p\) & \(\to\)i 2--6
\end{tabular}
\end{minipage}\hfill
\begin{minipage}[t]{0.49\textwidth}
\textbf{(\(\Leftarrow\)) clásica}\\[2pt]
\begin{tabular}{lll}
1 & \(\neg q\to\neg p\) & premisa\\ \cline{2-3}
2 & \multicolumn{1}{|l}{\(p\)} & sup.\ (\(\to\)i)\\
3 & \multicolumn{1}{|l}{\quad\(\neg q\)} & sup.\ (RAA)\\
4 & \multicolumn{1}{|l}{\quad\(\neg p\)} & \(\to\)e 1, 3\\
5 & \multicolumn{1}{|l}{\quad\(\bot\)} & \(\neg\)e 2, 4\\ \cline{3-3}
6 & \multicolumn{1}{|l}{\(q\)} & RAA 3--5\\ \cline{2-3}
7 & \(p\to q\) & \(\to\)i 2--6
\end{tabular}
\end{minipage}

\medskip
Las dos derivaciones tienen la misma forma. La diferencia está en la línea 6:
\(\neg\)i concluye la \emph{negación} de lo supuesto, mientras que RAA
concluye la fórmula \emph{sin} la negación supuesta. Por eso la primera
dirección vale en NJ y la segunda no.
\end{ejemplotesis}

\begin{nota}[\(\neg\)i no es RAA (Nota 7.60)]
\[
\frac{\Gamma,A\vdash\bot}{\Gamma\vdash\neg A}\ \neg\mathrm I
\qquad\qquad
\frac{\Gamma,\neg A\vdash\bot}{\Gamma\vdash A}\ \mathrm{RAA}
\]
Las dos «derivan una contradicción», pero solo la segunda es clásica. Con
\(\neg A:=A\to\bot\), la primera es un caso de \(\to\)I. La segunda equivale a
la eliminación de la doble negación. Es la pregunta más probable sobre lógica
clásica e intuicionista.
\end{nota}

%---------------------------------------------------------------------
\begin{esquema}{Esquema 15. Solidez: inducción sobre la prueba y análisis de la última regla}
\forma Para \(\Gamma\vdash\psi\Rightarrow\Gamma\vDash\psi\): inducción sobre
la derivación (Esquema 9). Para cada regla se toma una valuación \(v\) que
satisface \(\Gamma\) y se comprueba que hace verdadera la conclusión, usando
las hipótesis inductivas de las premisas.

\porque Cada regla \emph{preserva la verdad}: si las premisas son verdaderas
bajo \(v\), la conclusión también lo es.

\donde Teorema 9.13 (esbozo; detalles en Huth y Ryan, Teorema 1.35).
\end{esquema}

\begin{ejemplotesis}{Dos casos del Teorema 9.13}
\textbf{\(\to\)i.} La premisa es \(\Gamma,\varphi\vdash\chi\), así que la
hipótesis inductiva da \(\Gamma,\varphi\vDash\chi\). Sea \(v\vDash\Gamma\). Si
\(v(\varphi)=F\), entonces \(v(\varphi\to\chi)=T\). Si \(v(\varphi)=T\), la
hipótesis da \(v(\chi)=T\). En ambos casos \(v(\varphi\to\chi)=T\).

\textbf{RAA.} La premisa es \(\Gamma,\neg\varphi\vdash\bot\), así que la
hipótesis da \(\Gamma,\neg\varphi\vDash\bot\). Si fuera \(v(\varphi)=F\),
entonces \(v\) satisfaría \(\Gamma\cup\{\neg\varphi\}\) y la hipótesis daría
\(v(\bot)=T\), lo cual es imposible. Luego \(v(\varphi)=T\).
\end{ejemplotesis}

\begin{trampa}
El caso RAA es el único, junto con LEM, que usa la \emph{bivalencia}: el paso
de «no es \(F\)» a «es \(T\)». En una semántica no bivalente (por ejemplo, la
de Kripke) RAA deja de ser sólida, mientras que las reglas del núcleo
intuicionista siguen siéndolo. Esto explica por qué NJ no prueba \(p\vee\neg p\).
\end{trampa}

%---------------------------------------------------------------------
\begin{esquema}{Esquema 16. Reforzar el enunciado para que la inducción pase (Kalmár)}
\forma Si una inducción se atasca porque la hipótesis inductiva es demasiado
débil, se prueba un enunciado \emph{más fuerte} que la incluya. En el lema de
Kalmár, en lugar de «si \(v(\eta)=T\), los literales de \(v\) derivan
\(\eta\)», se prueba el enunciado \textbf{bilateral}:
\[
\Gamma_v\vdash\eta\ \text{ si } v(\eta)=T,
\qquad
\Gamma_v\vdash\neg\eta\ \text{ si } v(\eta)=F,
\]
donde \(\Gamma_v\) contiene, para cada átomo \(p\) de \(\eta\), el literal
\(p\) si \(v(p)=T\) y \(\neg p\) si \(v(p)=F\).

\porque En los casos \(\neg\) y \(\to\) la verdad de la fórmula depende de que
alguna subfórmula sea \emph{falsa}, y para eso hace falta una hipótesis
inductiva que hable también de las fórmulas falsas.

\donde Teorema 9.14 (esbozo, Huth y Ryan §1.4.4) y Ejemplo 9.17 (\(p\to p\)).
El Capítulo 11 lo formaliza.
\end{esquema}

\begin{ejemplotesis}{Los casos que obligan al enunciado bilateral (desarrollo del esbozo)}
\textbf{\(\eta\equiv\neg\varphi\).}
\begin{itemize}
  \item Si \(v(\neg\varphi)=T\), entonces \(v(\varphi)=F\), y la hipótesis
    \emph{sobre fórmulas falsas} da \(\Gamma_v\vdash\neg\varphi\). Con un
    enunciado unilateral no habría nada que usar.
  \item Si \(v(\neg\varphi)=F\), entonces \(v(\varphi)=T\), y la hipótesis da
    \(\Gamma_v\vdash\varphi\); de ahí \(\Gamma_v\vdash\neg\neg\varphi\), que es
    válido intuicionistamente.
\end{itemize}
\textbf{\(\eta\equiv\varphi\to\psi\).}
\begin{itemize}
  \item \(v(\varphi)=F\): de \(\Gamma_v\vdash\neg\varphi\) se obtiene
    \(\Gamma_v\vdash\varphi\to\psi\) (suponer \(\varphi\), contradicción,
    \(\bot\)e).
  \item \(v(\psi)=T\): de \(\Gamma_v\vdash\psi\) se obtiene
    \(\Gamma_v\vdash\varphi\to\psi\) (\(\to\)i sin usar la suposición).
  \item \(v(\varphi)=T\), \(v(\psi)=F\): de \(\Gamma_v\vdash\varphi\) y
    \(\Gamma_v\vdash\neg\psi\) se obtiene \(\Gamma_v\vdash\neg(\varphi\to\psi)\)
    (suponer \(\varphi\to\psi\), modus ponens, contradicción).
\end{itemize}
\textbf{Ensamblaje.} Si \(\eta\) es tautología, el primer caso ocurre en las
\(2^k\) filas. Los átomos se descargan uno a uno: de \(\Delta,p\vdash\eta\) y
\(\Delta,\neg p\vdash\eta\) se obtiene \(\Delta\vdash\eta\). Este es el único
paso clásico: usa LEM y \(\vee\)e, o RAA en el fragmento \(\{\neg,\to\}\).
\end{ejemplotesis}

\begin{ejemplotesis}{Ejemplo 9.17: completitud como construcción para \(p\to p\)}
Para cada fila se da una derivación:
\begin{center}
\AxiomC{}\RightLabel{\scriptsize hip.}\UnaryInfC{\(p\vdash p\)}
\RightLabel{\scriptsize\(\to\)I}\UnaryInfC{\(p\vdash p\to p\)}
\DisplayProof
\qquad\qquad
\AxiomC{}\RightLabel{\scriptsize hip.}\UnaryInfC{\(\neg p,p\vdash p\)}
\RightLabel{\scriptsize\(\to\)I}\UnaryInfC{\(\neg p\vdash p\to p\)}
\DisplayProof
\end{center}
El análisis por casos clásico sobre \(p\) las combina en \(\vdash p\to p\). El
ejemplo es pequeño, pero muestra el programa completo: una derivación por
fila de la tabla y luego la descarga de los átomos.
\end{ejemplotesis}

%---------------------------------------------------------------------
\begin{esquema}{Esquema 17. Argumentos de imposibilidad}
\forma Para probar que algo \emph{no} existe (un tipo, una derivación), se
supone que existe y se deriva una restricción imposible.

\porque Es la introducción de la negación en el metalenguaje: suponer y llegar
a una contradicción.

\donde Ejemplo 7.49(3): \(x\,x\) no es tipable. Ejemplo 9.20: tautologías
clásicas que no son válidas bajo BHK.
\end{esquema}

\begin{ejemplotesis}{Ejemplo 7.49(3): \(x\,x\) no es tipable en \(\lambda\!\to\)}
Si \(x\,x:\tau\), la regla (aplicación) exige \(x:\sigma\to\tau\) (como
función) y \(x:\sigma\) (como argumento). Como \(x\) tiene un solo tipo en el
contexto, \(\sigma\equiv\sigma\to\tau\). Eso es imposible: el tipo de la
derecha tiene estrictamente más símbolos que el de la izquierda, y dos tipos
idénticos tienen el mismo tamaño.
\end{ejemplotesis}

\begin{trampa}[Qué demuestra el Ejemplo 9.20 y qué no]
El ejemplo argumenta en la lectura informal BHK: «no existe un método
universal». Es un argumento de motivación, no una prueba de no derivabilidad
en NJ. La prueba rigurosa de que \(p\vee\neg p\) no es un teorema
intuicionista requiere un contramodelo de Kripke y la solidez de NJ respecto
de esa semántica. La tesis lo remite a Sørensen y Urzyczyn (Nota 9.18) y no lo
formaliza. Conviene decirlo así si se pregunta.
\end{trampa}

\subsection*{Parte E. Pruebas como programas}

%---------------------------------------------------------------------
\begin{esquema}{Esquema 18. Leer una prueba como un programa (Curry--Howard)}
\forma A cada derivación se le asocia un término, y viceversa:
\begin{center}
\begin{tabular}{@{}ll@{}}
\toprule
\textbf{Lógica} & \textbf{Cómputo}\\
\midrule
suposición \(A\) (abre caja) & variable \(x:A\)\\
\(\to\)i (cierra caja) & abstracción \(\lambda x.M\)\\
\(\to\)e (modus ponens) & aplicación \(M\,N\)\\
\(\wedge\)i / \(\wedge\)e\(_i\) & par \(\langle M,N\rangle\) / proyección \(\pi_i\)\\
\(\vee\)i\(_i\) / \(\vee\)e & inyección \(\mathrm{in}_i\) / \(\operatorname{case}\)\\
\(\forall\)i / \(\forall\)e & abstracción dependiente / aplicación a un término\\
\bottomrule
\end{tabular}
\end{center}

\porque Al borrar los términos de las reglas de tipado (Figura de Church) y
cambiar \(\times,+\) por \(\wedge,\vee\), se obtienen exactamente las reglas de
Gentzen. El Teorema 9.22 (citado de Sørensen y Urzyczyn) lo enuncia para el
fragmento implicacional.

\plantilla Para explicar una prueba, escribir su término y leerlo como
algoritmo: «recibe una prueba de \dots\ y devuelve \dots».

\donde Ejemplo 9.24 (\emph{swap}), Definición 9.19 (BHK), Ejemplo 7.52 y
§8.3.3.
\end{esquema}

\begin{ejemplotesis}{Ejemplo 9.24: \((B\wedge A)\to(A\wedge B)\) es \(\lambda z.\langle\pi_2 z,\pi_1 z\rangle\)}
\begin{prooftree}
\AxiomC{\([z:B\times A]\)}\UnaryInfC{\(\pi_2 z:A\)}
\AxiomC{\([z:B\times A]\)}\UnaryInfC{\(\pi_1 z:B\)}
\BinaryInfC{\(\langle\pi_2 z,\pi_1 z\rangle:A\times B\)}
\RightLabel{\scriptsize\(\to\)I\(^z\)}
\UnaryInfC{\(\lambda z.\langle\pi_2 z,\pi_1 z\rangle:(B\times A)\to(A\times B)\)}
\end{prooftree}
Al borrar los términos queda la prueba lógica, y al conservarlos queda una
función que intercambia los componentes de un par.
\end{ejemplotesis}

\begin{trampa}
\begin{itemize}
  \item La correspondencia no asigna un término \emph{único} a cada
    secuente: \(\varphi\to(\varphi\to\varphi)\) tiene los habitantes
    \(\lambda x.\lambda y.x\) y \(\lambda x.\lambda y.y\) (§9.5, tras la
    Nota 9.23). Es
    biyectiva entre derivaciones con hipótesis etiquetadas y términos.
  \item Es una correspondencia para la lógica \emph{intuicionista}. El cálculo
    del Capítulo 11 es clásico (tiene RAA); allí la derivabilidad se
    representa con un predicado inductivo, lo cual no la vuelve intuicionista
    (Nota 9.18).
\end{itemize}
\end{trampa}

%---------------------------------------------------------------------
\begin{esquema}{Esquema 19. Normalización de pruebas es evaluación de programas}
\forma Un \emph{desvío} es una introducción seguida inmediatamente de la
eliminación del mismo conectivo. Eliminarlo es un paso de reducción:
\[
(\lambda x.N)\,M\to N[x:=M],\qquad
\pi_i\langle M_1,M_2\rangle\to M_i,
\]
\[
\operatorname{case}(\mathrm{in}_1(N);x.K;y.L)\to K[x:=N],\qquad
\operatorname{case}(\mathrm{in}_2(N);x.K;y.L)\to L[y:=N].
\]

\porque En la prueba con desvío, la suposición \(A\) se usa en algún lugar de
la prueba de \(B\). Reemplazar cada uso por la prueba de \(A\) que se tenía da
una prueba directa de \(B\). Ese reemplazo es literalmente la sustitución
\(N[x:=M]\) de la reducción \(\beta\).

\donde Definición 9.25 y Figura «Simplificación de pruebas frente a evaluación
de programas» (§9.5.2).
\end{esquema}

\begin{nota}[Por qué importa para la tesis]
Esta lectura explica por qué los metateoremas de la reducción del
Capítulo 7 (confluencia y normalización) son también metateoremas sobre
\emph{pruebas}. Por ejemplo, la normalización fuerte de los términos tipados
(Teorema 8.27) dice que toda prueba puede llevarse a una forma sin desvíos.
\end{nota}

\subsection*{Parte F. Tipos inductivos (Capítulo 10)}

%---------------------------------------------------------------------
\begin{esquema}{Esquema 20. Recursión e inducción a partir del eliminador}\label{esq:elim}
\forma Para un tipo inductivo con constructores \(0:\mathbb N\) y
\(S:\mathbb N\to\mathbb N\):
\begin{itemize}
  \item \textbf{Definir una función} \(g:\mathbb N\to C\) es dar un valor base
    \(x:C\) y una regla de paso \(f:\mathbb N\to C\to C\). El eliminador no
    dependiente da \(g(n):=R(C,x,f,n)\), con
    \[
    R(C,x,f,0)\to_\iota x,\qquad R(C,x,f,S\,n)\to_\iota f\,n\,R(C,x,f,n).
    \]
  \item \textbf{Probar \(\forall n.\,P\,n\)} es dar \(p_0:P\,0\) y
    \(p_S:\forall u.\,P\,u\to P\,(S\,u)\). El eliminador dependiente tiene
    exactamente el tipo del principio de inducción:
    \[
    \forall P:\mathbb N\to\mathrm{Prop},\ P\,0\to(\forall u,\ P\,u\to P\,(S\,u))\to\forall n,\ P\,n.
    \]
\end{itemize}

\porque Los constructores son las únicas formas de construir un habitante, y
el eliminador es la única forma de usarlo (Nota 10.6). La \(\iota\)-reducción
fija cómo se calcula el eliminador frente a cada constructor.

\plantilla «Se define \(g\) por recursión: \(g(0)=x\) y \(g(n+1)=f(n,g(n))\)».
«Por inducción sobre \(n\): caso \(0\), \dots; caso \(S\,u\), suponiendo
\(P\,u\), \dots». Son las frases de siempre; el Capítulo 10 explica de dónde
salen.

\donde §10.3.1 (recursor y Nota 10.1), Ejemplo 10.4, §10.3.4
(eliminación dependiente y no dependiente) y §10.3.5 (\(\iota\)).
\end{esquema}

\begin{ejemplotesis}{Nota 10.1: el factorial con el recursor}
\(\mathrm{fact}(n):=R(\mathbb N,\,1,\,\lambda n\,r.\,(S\,n)\cdot r,\,n)\). Al
calcular, cada paso es una \(\iota\)-reducción seguida de una \(\beta\):
\[
\begin{aligned}
R(\mathbb N,1,f,S(S\,0))
&\to_\iota f\,(S\,0)\,R(\mathbb N,1,f,S\,0)
\ \to_\beta\ 2\cdot R(\mathbb N,1,f,S\,0)\\
&\to_\iota 2\cdot\big(f\,0\,R(\mathbb N,1,f,0)\big)
\ \to_\beta\ 2\cdot\big(1\cdot R(\mathbb N,1,f,0)\big)
\ \to_\iota\ 2\cdot(1\cdot 1)=2.
\end{aligned}
\]
\end{ejemplotesis}

\begin{nota}[Unificación: todas las inducciones son eliminadores]
La inducción estructural sobre \(\Lambda\) (Esquema 1) es el eliminador
dependiente del tipo inductivo de los términos: un caso por constructor y una
hipótesis inductiva por cada argumento recursivo (\(P\) y \(Q\) en \(PQ\);
\(P\) en \(\lambda x.P\)). La inducción sobre derivaciones (Esquema 9) es el
eliminador de la \emph{familia inductiva} \(\Gamma\vdash M:\sigma\), indexada
por \(\Gamma\), \(M\) y \(\sigma\); por eso el contexto puede variar entre la
premisa y la conclusión. El Capítulo 10 muestra que los esquemas de los
capítulos 7 a 9 son instancias de una sola regla del CIC.
\end{nota}

%---------------------------------------------------------------------
\begin{esquema}{Esquema 21. Codificar por el uso (codificaciones impredicativas)}
\forma En el CoC puro, un dato o una proposición se define por \emph{cómo se
usa}:
\[
\begin{aligned}
\mathrm{Nat}&:=\forall A.\ A\to(A\to A)\to A,&\quad n&=\lambda A\,x\,f.\,f^{\,n}x,\\
A\wedge B&:=\forall C.\ (A\to B\to C)\to C,&\quad a=_L b&:=\forall P.\ P\,a\to P\,b.
\end{aligned}
\]
Las pruebas sobre estas codificaciones son cálculos \(\beta\) o
instanciaciones bien elegidas.

\porque El producto dependiente \(\Pi\) permite cuantificar sobre todas las
proposiciones, incluida la que se está definiendo (impredicatividad).

\donde §10.2.2 (Tabla de codificaciones de Coquand y Huet) y §10.2.3
(numerales, predecesor e igualdad de Leibniz).
\end{esquema}

\begin{ejemplotesis}{Dos cálculos típicos}
\textbf{Sucesor.} Con \(S:=\lambda n\,A\,x\,f.\,f\,(n\,A\,x\,f)\),
\[
S\,1\ \bred\ \lambda A\,x\,f.\,f\,(1\,A\,x\,f)
\ \bmred\ \lambda A\,x\,f.\,f\,(f\,x)\ \equiv\ 2.
\]
\textbf{Simetría de la igualdad de Leibniz (el truco del predicado).} De
\(h:\forall P.\ P\,a\to P\,b\) se quiere \(\forall P'.\ P'\,b\to P'\,a\). Se
instancia \(h\) en \(P(z):=P'(z)\to P'(a)\):
\[
h\,P:\ (P'\,a\to P'\,a)\to(P'\,b\to P'\,a),
\qquad\text{luego}\qquad
h\,P\,(\lambda t.\,t):\ P'\,b\to P'\,a.
\]
Es la instanciación que la tesis menciona en §10.2.3 para mostrar que las dos
direcciones de la definición son equivalentes. El patrón general consiste en
elegir el predicado de modo que el caso conocido sea trivial (aquí, la
identidad).
\end{ejemplotesis}

\begin{trampa}[El límite que motiva el CIC]
La codificación da \emph{iteración}, pero no permite derivar el principio de
inducción dependiente \(\forall P.\,P\,0\to(\forall n.\,P\,n\to P\,(S\,n))\to\forall n.\,P\,n\).
Es un resultado de no derivabilidad de Geuvers, citado y no demostrado en la
tesis. Además, el predecesor cuesta \(n\) iteraciones con la construcción por
pares. Estas dos limitaciones justifican pasar a definiciones inductivas
primitivas.
\end{trampa}

%---------------------------------------------------------------------
\begin{esquema}{Esquema 22. Por qué la positividad estricta (complemento)}
\forma Las ocurrencias recursivas de \(X\) en los constructores no pueden estar
a la izquierda de una flecha (§10.3.2). Para justificar la restricción se
muestra que, sin ella, se obtiene un habitante de \(\bot\).

\porque Este es el argumento que la tesis resume en la frase «se podría
construir un punto fijo y obtener un habitante de \(\bot\)». Supóngase un tipo
\(X\) con un único constructor \(c:(X\to\bot)\to X\). El análisis de casos
daría \(u:X\to(X\to\bot)\) con \(u\,(c\,g)\to g\). Sea
\[
d:=\lambda x{:}X.\ u\,x\,x\ :\ X\to\bot.
\]
Entonces \(d\,(c\,d):\bot\). Es el argumento diagonal de Cantor y Russell en
forma de programa: \(d\,(c\,d)\to u\,(c\,d)\,(c\,d)\to d\,(c\,d)\to\cdots\)
no termina, y la lógica se vuelve inconsistente.

\donde §10.3.2 (tipos estrictamente positivos).
\end{esquema}

\subsection*{Parte G. Resultados que la tesis cita}

%---------------------------------------------------------------------
\begin{esquema}{Esquema 23. Presentar un resultado citado}
\forma Para un teorema que la tesis no demuestra, conviene decir tres cosas:
qué afirma, de dónde se toma y para qué se usa en la tesis.

\porque Distinguir lo demostrado de lo citado es parte de la honestidad del
trabajo, y es lo primero que verifica un examinador.

\plantilla
\begin{center}
\begin{tabularx}{\linewidth}{@{}p{0.2\linewidth}X@{}}
\toprule
\textbf{Resultado} & \textbf{Fuente y papel en la tesis}\\
\midrule
Teorema 7.37 & Barendregt, Teorema 3.2.8; esbozo de Tait--Martin-Löf. Se usa en
los Corolarios 7.43 y 7.44.\\
Lemas 7.40, 7.41 & Enunciados «sin demostración» dentro del esbozo de 7.37.\\
Teorema 7.59 & Barendregt (2013). Decidibilidad de tipabilidad, chequeo y
habitación en \(\lambda\!\to\).\\
Teorema 7.63 & Barendregt (1992). Church--Rosser en \(\lambda\!\to\).\\
Teorema 8.27 & Barendregt (1991); Sørensen y Urzyczyn. Confluencia,
preservación del tipo, unicidad de tipos y normalización fuerte en el cubo.\\
Teorema 9.22 & Sørensen y Urzyczyn. Isomorfismo de Curry--Howard.\\
Teoremas 9.13, 9.14 & Huth y Ryan (Teorema 1.35, §1.4.4). Esbozos; la
completitud se formaliza en el Capítulo 11.\\
§10.2.3 & Geuvers (2001). La inducción no se deriva de la codificación de
Church.\\
\bottomrule
\end{tabularx}
\end{center}
\end{esquema}

\begin{nota}[Una distinción que la tesis hace explícita (§8.4.3)]
La \emph{comprobación} de tipos \(\Gamma\overset{?}{\vdash}M:A\) es decidible
en el cubo, pero la \emph{habitación} \(\exists M.\ \Gamma\vdash M:A\) es
indecidible en \(\lambda P\) y en \(\lambda C\). Dicho de otro modo,
comprobar una prueba dada es mecánico, mientras que encontrar una prueba no lo
es. Esta es la razón de ser de un asistente de pruebas: el humano propone y el
núcleo comprueba.
\end{nota}

%=====================================================================
\section{Lecturas guiadas}
%=====================================================================

\subsection{El Lema de sustitución (Lema 7.23), de principio a fin}

\textbf{Enunciado.} Si \(x\not\equiv y\) y \(x\notin\FV(L)\), entonces
\(M[x:=N][y:=L]\equiv M[y:=L][x:=N[y:=L]]\).

\textbf{Qué dice en palabras.} Da igual sustituir primero \(x\) por \(N\) y
luego \(y\) por \(L\), que sustituir primero \(y\) por \(L\) y luego \(x\) por
\(N[y:=L]\). La segunda forma «compensa» que \(N\) también deba recibir la
sustitución de \(y\).

\textbf{Estructura de la prueba.}
\begin{enumerate}
  \item \emph{Esquema 1}, inducción sobre \(M\). Los términos \(N\), \(L\) y
    las variables \(x,y\) se mantienen fijos.
  \item \emph{Caso variable: Esquema 2}, tres subcasos. Cada hipótesis del
    lema se usa en uno: \(x\not\equiv y\) en \(z\equiv x\) y
    \(x\notin\FV(L)\) en \(z\equiv y\).
  \item \emph{Caso aplicación: Esquema 3}, una cadena de tres igualdades:
    definición, hipótesis de inducción y definición leída al revés.
  \item \emph{Caso abstracción: Esquemas 4 y 3}. Se renombra la variable
    ligada a una \(z\) fresca, la sustitución entra bajo \(\lambda\) en ambos
    lados y se aplica la hipótesis de inducción.
\end{enumerate}

\textbf{Qué responder si preguntan.}
\begin{itemize}
  \item \emph{¿Por qué hace falta \(x\notin\FV(L)\)?} Contraejemplo con
    \(M\equiv y\), \(L\equiv x\): el lado izquierdo es \(y[x:=N][y:=x]\equiv x\)
    y el derecho es \(x[x:=N[y:=x]]\equiv N[y:=x]\), que en general no es
    \(x\).
  \item \emph{¿Por qué es «módulo \(\alpha\)»?} Porque la cláusula (3) de la
    sustitución elige una variable fresca que no es única.
  \item \emph{¿Para qué sirve?} Es el ingrediente del Lema 7.40 (sustitución
    para \(\pr\)) y, por tanto, del esbozo de Church--Rosser.
\end{itemize}

\subsection{De Church--Rosser a la unicidad de la forma normal}

La cadena lógica de §7.2.8 es corta y conviene poder recitarla:
\[
\underbrace{\text{Lemas 7.40, 7.41}}_{\pr\ \text{tiene diamante}}
\ \Longrightarrow\
\underbrace{\text{Teorema 7.37}}_{\bmred\ \text{confluye}}
\ \Longrightarrow\
\underbrace{\text{Corolario 7.43}}_{\beq\ \text{tiene reducto común}}
\ \Longrightarrow\
\underbrace{\text{Corolario 7.44}}_{\text{forma normal única}}
\]
\begin{itemize}
  \item Primera flecha: Esquema 8 (embaldosado y \(\bred\subseteq\pr\subseteq\bmred\)).
  \item Segunda: Esquema 6 (inducción sobre la longitud de la conversión).
  \item Tercera: Esquema 7 (las formas normales solo se reducen a sí mismas).
\end{itemize}
\textbf{Consecuencia que conviene mencionar:} \(\beq\) no identifica formas
normales distintas. Por ejemplo, \(\lambda x.\lambda y.x\not\beq\lambda x.\lambda y.y\),
así que la teoría \(\beta\) es consistente en el sentido de no igualar todos los
términos.

\subsection{Una misma frase, cuatro lenguajes}

La frase «si \(Q\) vale para todos los pares, vale para la diagonal» aparece
en la tesis (§8.3.3) y permite mostrar en un solo ejemplo cómo encajan los
capítulos:
\begin{center}
\renewcommand{\arraystretch}{1.3}
\begin{tabularx}{\linewidth}{@{}p{0.27\linewidth}X@{}}
\toprule
\textbf{Lenguaje} & \textbf{La misma prueba}\\
\midrule
Matemática informal & Sea \(u\) arbitrario; por hipótesis, \(Q(u,u)\).\\
Deducción natural & Suposición \(\forall x\forall y.Q(x,y)\); \(\forall\)e dos veces; \(\forall\)i; \(\to\)i.\\
Derivación en \(\lambda P\) & Banderas \(z\), \(u\); (aplicación) dos veces; (abstracción) dos veces.\\
Término & \(\lambda z.\,\lambda u{:}S.\ z\,u\,u\)\\
\bottomrule
\end{tabularx}
\end{center}

%=====================================================================
\section{Cómo explicarlo en la defensa}
%=====================================================================

\subsection{Receta de cinco pasos para cualquier demostración}

\begin{enumerate}
  \item \textbf{Decir el enunciado en palabras}, sin símbolos.
  \item \textbf{Nombrar el esquema}: «es una inducción sobre la derivación»,
    «es una cadena de reducciones», «se construye la derivación de abajo
    hacia arriba».
  \item \textbf{Nombrar el objeto sobre el que se induce o se analiza}, y por
    qué ese y no otro.
  \item \textbf{Ir al caso interesante} (no al rutinario): la abstracción, con
    su variable fresca; el último eslabón de la cadena; RAA en la solidez.
  \item \textbf{Decir qué hipótesis se usó dónde}. Si una hipótesis no se
    usa, o es superflua o hay un error.
\end{enumerate}

\subsection{Diccionario de frases}

{\small
\begin{longtable}{@{}p{0.45\textwidth}p{0.5\textwidth}@{}}
\toprule
\textbf{Frase en la tesis} & \textbf{Qué significa formalmente}\\
\midrule
\endhead
«Por inducción estructural sobre \(M\)» & Un caso por cada regla de formación
de términos; hipótesis inductiva para los subtérminos inmediatos.\\
«Por inducción sobre la derivación» & Un caso por cada regla de tipado o de
deducción; hipótesis inductiva para las premisas.\\
«Se puede renombrar \(z\)» & Se elige un representante de la clase de
\(\alpha\)-equivalencia con \(z\) fresca.\\
«Existe una secuencia de reducción» & Testigo del existencial que define
\(\bmred\) o \(\beq\).\\
«Al concatenar ambas» & Transitividad: la cadena de longitud \(r+s\).\\
«Invirtiendo la secuencia» & Simetría de \(\beq\), posible por la disyunción en
cada eslabón.\\
«Se aplica la regla (abstracción)» & Se descarga la última declaración del
contexto (se cierra la bandera o la caja).\\
«Suposición (para \(\to\)i)» & Se abre una caja; su alcance termina al
aplicar \(\to\)i.\\
«Por RAA» & Se supuso \(\neg\varphi\), se llegó a \(\bot\) y se concluye
\(\varphi\). Es el paso clásico.\\
«El término habita el tipo» & Existe una derivación de \(\vdash M:A\); bajo
PAT, \(A\) es demostrable.\\
«Se reduce por \(\iota\)» & El eliminador aplicado a un constructor se
calcula con la rama de ese constructor.\\
\bottomrule
\end{longtable}
}

\subsection{Preguntas frecuentes}

\begin{description}[style=nextline, leftmargin=0pt, itemsep=6pt]
\item[¿Qué diferencia hay entre inducción estructural e inducción sobre derivaciones?]
La primera recorre las reglas que construyen \emph{términos}; la segunda, las
reglas que construyen \emph{juicios}. En \(\lambda\!\to\) coinciden porque el
tipado está dirigido por la sintaxis. Ambas son instancias del eliminador de
un tipo inductivo (Capítulo 10).

\item[¿Church--Rosser implica que todo término tiene forma normal?]
No. Solo implica que, si existe, es única (Corolario 7.44). \(\Omega\) no tiene
forma normal. Para términos tipados, la existencia la da la normalización
fuerte (Teorema 8.27, citado).

\item[¿Por qué se necesita la reducción paralela?]
Porque \(\bred\) no tiene el diamante en un paso: un redex duplicado obliga a
dar dos pasos en un lado. \(\pr\) contrae todas las copias a la vez
(Esquema~8).

\item[¿Por qué el Lema 7.23 es «módulo \(\alpha\)»?]
Porque la sustitución en una abstracción elige una variable fresca, y esa
elección no es única. La tesis identifica términos \(\alpha\)-equivalentes
desde §7.2.4.

\item[¿En qué se diferencian \(\neg\)i y RAA?]
\(\neg\)i concluye \(\neg A\) de \(A\vdash\bot\) y es intuicionista. RAA
concluye \(A\) de \(\neg A\vdash\bot\) y es clásica. En la Proposición 9.9
la primera dirección usa \(\neg\)i y la segunda RAA.

\item[¿Por qué el lema de Kalmár es bilateral?]
Porque en los casos \(\neg\) y \(\to\) la verdad de la fórmula depende de
subfórmulas falsas; la hipótesis inductiva tiene que dar
\(\Gamma_v\vdash\neg\varphi\) cuando \(v(\varphi)=F\).

\item[¿Dónde entra la lógica clásica en la completitud?]
Solo al descargar los átomos: de \(\Delta,p\vdash\eta\) y
\(\Delta,\neg p\vdash\eta\) a \(\Delta\vdash\eta\). Eso requiere LEM, o RAA en
el fragmento \(\{\neg,\to\}\). Por eso la completitud respecto de las tablas
de verdad falla para NJ.

\item[¿Qué aporta el CIC frente al CoC?]
El principio de inducción. En el CoC los naturales de Church dan iteración,
pero la inducción dependiente no es derivable (Geuvers). En el CIC, declarar
un tipo inductivo genera su eliminador, cuyo tipo es la inducción.

\item[¿Por qué se exige positividad estricta?]
Porque un constructor como \(c:(X\to\bot)\to X\) permite el término diagonal
\(d\,(c\,d):\bot\) (Esquema 22), y la lógica se vuelve inconsistente.

\item[¿Comprobar una prueba es lo mismo que encontrarla?]
No. El chequeo de tipos es decidible; la habitación no lo es en \(\lambda P\)
ni en \(\lambda C\) (§8.4.3). El asistente comprueba; el humano (o la
búsqueda) propone.

\item[¿Qué partes del capítulo 7 están demostradas completas?]
Los Lemas 7.6, 7.23, 7.31, 7.32 y 7.57 y los Corolarios 7.43 y 7.44.
Church--Rosser (7.37) está esbozado; sus lemas auxiliares (7.40 y 7.41) solo
se enuncian; los Teoremas 7.59 y 7.63 se citan.
\end{description}

%=====================================================================
\appendix
\section{Hoja de referencia: reglas}
%=====================================================================

\subsection*{\(\lambda\!\to\) (§7.3.3)}
\[
\frac{x:\sigma\in\Gamma}{\Gamma\vdash x:\sigma}\ \text{(var)}
\qquad
\frac{\Gamma\vdash M:\sigma\to\tau\quad\Gamma\vdash N:\sigma}{\Gamma\vdash MN:\tau}\ \text{(apl)}
\qquad
\frac{\Gamma,x:\sigma\vdash M:\tau}{\Gamma\vdash\lambda x{:}\sigma.M:\sigma\to\tau}\ \text{(abs)}
\]

\subsection*{\(\lambda2\) (Definición 8.4)}
\[
\frac{\Gamma,\alpha:*\vdash M:A}{\Gamma\vdash\Lambda\alpha{:}*.M:\Pi\alpha{:}*.A}\ \text{(abs\(_2\))}
\qquad
\frac{\Gamma\vdash M:\Pi\alpha{:}*.A\quad\Gamma\vdash B:*}{\Gamma\vdash M\,B:A[\alpha:=B]}\ \text{(apl\(_2\))}
\]

\subsection*{\(\lambda\underline{\omega}\) (§8.2)}
\[
\varnothing\vdash *:\square\ \text{(sort)}
\qquad
\frac{\Gamma\vdash A:s}{\Gamma,x:A\vdash x:A}\ \text{(var)}
\qquad
\frac{\Gamma\vdash A:B\quad\Gamma\vdash C:s}{\Gamma,x:C\vdash A:B}\ \text{(deb)}
\]
\[
\frac{\Gamma\vdash A:s\quad\Gamma\vdash B:s}{\Gamma\vdash A\to B:s}\ \text{(form)}
\qquad
\frac{\Gamma\vdash M:A\to B\quad\Gamma\vdash N:A}{\Gamma\vdash MN:B}\ \text{(apl)}
\]
\[
\frac{\Gamma,x:A\vdash M:B\quad\Gamma\vdash A\to B:s}{\Gamma\vdash\lambda x{:}A.M:A\to B}\ \text{(abs)}
\qquad
\frac{\Gamma\vdash A:B\quad\Gamma\vdash B':s}{\Gamma\vdash A:B'}\ \text{(conv, si \(B\beq B'\))}
\]

\subsection*{Cubo \(\lambda\) (§8.4.2)}
\[
\frac{\Gamma\vdash A:s_1\quad\Gamma,x:A\vdash B:s_2}{\Gamma\vdash\Pi x{:}A.B:s_2}
\qquad
\begin{array}{ll}
(*,*) & \lambda\!\to\ \text{(términos de términos)}\\
(\square,*) & \lambda2\ \text{(términos de tipos)}\\
(*,\square) & \lambda P\ \text{(tipos de términos)}\\
(\square,\square) & \lambda\underline{\omega}\ \text{(tipos de tipos)}
\end{array}
\]
\(\lambda C\) admite los cuatro pares.

\subsection*{Deducción natural (§9.1) y NJ (§9.4.2)}
\[
\frac{\varphi\quad\psi}{\varphi\wedge\psi}\wedge\mathrm i
\quad
\frac{\varphi\wedge\psi}{\varphi}\wedge\mathrm e_1
\quad
\frac{\varphi}{\varphi\vee\psi}\vee\mathrm i_1
\quad
\frac{\varphi\vee\psi\quad[\varphi]\cdots\chi\quad[\psi]\cdots\chi}{\chi}\vee\mathrm e
\quad
\frac{[\varphi]\cdots\psi}{\varphi\to\psi}\to\mathrm i
\quad
\frac{\varphi\quad\varphi\to\psi}{\psi}\to\mathrm e
\]
\[
\frac{\varphi\quad\neg\varphi}{\bot}\neg\mathrm e
\quad
\frac{[\varphi]\cdots\bot}{\neg\varphi}\neg\mathrm i
\quad
\frac{\bot}{\varphi}\bot\mathrm e
\qquad\Big|\qquad
\frac{[\neg\varphi]\cdots\bot}{\varphi}\mathrm{RAA}
\quad
\varphi\vee\neg\varphi\ \ \mathrm{LEM}
\]
NJ es el bloque de la izquierda (con \(\neg\varphi:=\varphi\to\bot\)); la lógica
clásica añade RAA o LEM.

\subsection*{CIC (§10.3)}
\[
\frac{\Gamma,X:A\vdash C_i:s}{\Gamma\vdash\Ind(X:A)\{C_1|\dots|C_n\}:A}
\qquad
\frac{\Gamma\vdash I:A}{\Gamma\vdash\Constr(i,I):C_i[X:=I]}
\]
\[
\Elim(\Constr(i,I)\,\vec m,\,Q)\{\vec f\}\ \to_\iota\ f_i\ \vec m
\]
donde \(f_i\) recibe, además de \(\vec m\), las llamadas recursivas sobre los argumentos
recursivos de \(\vec m\).

\end{document}
